% Sections 12.1 to 12.6, from lectures/L12/appendix/a_theory.md. Its unnumbered opening, "Varför
% derivata i elektroteknik?", is the introduction in chapter.tex.

\section{Förändringshastighet och ändringskvot}
\label{c12:sec:andringskvot}
\textbf{Ändringskvoten} (medelhastigheten) av $f(x)$ över intervallet $[x, x + h]$ är
\[
  \frac{f(x + h) - f(x)}{h}.
\]
Geometriskt är detta lutningen av \textbf{sekanten} genom punkterna $(x, f(x))$ och
$(x + h, f(x + h))$.

\section{Derivatans definition}
\label{c12:sec:definition}
Derivatan $f'(x)$ är gränsvärdet av ändringskvoten när $h \to 0$:
\[
  f'(x) = \lim_{h \to 0} \frac{f(x + h) - f(x)}{h}
\]
Geometriskt är $f'(x)$ \textbf{lutningen av tangenten} till kurvan $y = f(x)$ i punkten $x$.
Derivatan kallas också \textbf{momentan förändringshastighet}.

\section{Deriveringsregler för polynom}
\label{c12:sec:regler}

\begin{narrowtable}{@{}ll@{}}
  \thead{Funktion} & \thead{Derivata} \\
  \midrule
  $f(x) = C$ (konstant) & $f'(x) = 0$ \\
  $f(x) = x^n$ & $f'(x) = nx^{n-1}$ \\
  $f(x) = Cx^n$ & $f'(x) = Cnx^{n-1}$ \\
  $f(x) = u(x) + v(x)$ & $f'(x) = u'(x) + v'(x)$ \\
\end{narrowtable}

\begin{exempel}
\[
  f(x) = 3x^4 - 2x^2 + 5 \quad \Rightarrow \quad f'(x) = 12x^3 - 4x
\]
\end{exempel}

\section{Tangentens ekvation}
\label{c12:sec:tangent}
Tangenten till $y = f(x)$ i punkten $(x_0, f(x_0))$ har ekvationen
\[
  y = f'(x_0)(x - x_0) + f(x_0).
\]

\section{Extrempunkter -- maximum och minimum}
\label{c12:sec:extrempunkter}
En \textbf{stationär punkt} uppfyller $f'(x_0) = 0$. Om den är ett maximum eller ett minimum
avgörs med ett test med andra derivatan:
\begin{itemize}
  \item $f''(x_0) > 0$: \textbf{minimum} (kurvan är konkav uppåt),
  \item $f''(x_0) < 0$: \textbf{maximum} (kurvan är konkav nedåt).
\end{itemize}

\section{Typexempel}
\label{c12:sec:typexempel}

\begin{exempel}[Derivera polynomfunktioner]
Derivera
\begin{align*}
  \textbf{a)}\quad f(x) &= -2x^2 + 2x + 4, \\
  \textbf{b)}\quad f(x) &= 3x^3 - 6x^2 + \frac{3x}{4} - 5, \\
  \textbf{c)}\quad f(x) &= -x^4 + x^3 + \frac{2x^2}{3} - 3x + 2.
\end{align*}
Regeln för $Cx^n$ används term för term:
\begin{align*}
  \textbf{a)}\quad f'(x) &= -4x + 2 \\
  \textbf{b)}\quad f'(x) &= 9x^2 - 12x + \frac{3}{4} \\
  \textbf{c)}\quad f'(x) &= -4x^3 + 3x^2 + \frac{4x}{3} - 3
\end{align*}
\end{exempel}

\begin{exempel}[Analys av en parabel]
Funktionen är $f(x) = -x^2 + 6x - 5$.
\par\noindent\textbf{a)} Bestäm $f'(x)$.
\[
  f'(x) = -2x + 6
\]
\textbf{b)} Lös $f'(x) = 0$.
\[
  -2x + 6 = 0 \quad \Rightarrow \quad x = 3
\]
\textbf{c)} Avgör typ med $f''(x)$.
\[
  f''(x) = -2 < 0 \quad \Rightarrow \quad \text{maximum}
\]
\textbf{d)} Beräkna maxvärdet.
\[
  f(3) = -9 + 18 - 5 = 4
\]
Maximipunkten är $(3, 4)$.
\end{exempel}

\begin{exempel}[Strömförlopp i en RC-krets]
Strömmen approximeras av $i(t) = -0{,}4t^3 + 2{,}4t^2 - 3t + 1{,}2$ ampere, med tiden $t$ i
sekunder.
\par\noindent\textbf{a)} Beräkna $i'(t)$.
\[
  i'(t) = -1{,}2t^2 + 4{,}8t - 3
\]
\textbf{b)} Lös $i'(t) = 0$ med PQ-formeln.
\begin{gather*}
  t^2 - 4t + 2{,}5 = 0 \quad \Rightarrow \quad t = 2 \pm \sqrt{1{,}5} \\
  t_1 \approx 0{,}78\,\text{s}, \qquad t_2 \approx 3{,}22\,\text{s}
\end{gather*}
\textbf{c)} Andra derivatan är $i''(t) = -2{,}4t + 4{,}8$, och
\begin{alignat*}{2}
  i''(0{,}78) &\approx 2{,}9 > 0 &\qquad& \Rightarrow \quad \text{minimum vid } t_1, \\
  i''(3{,}22) &\approx -2{,}9 < 0 && \Rightarrow \quad \text{maximum vid } t_2.
\end{alignat*}
\textbf{d)} Strömmen i de stationära punkterna är
\[
  i(0{,}78) \approx 0{,}13\,\text{A} \ \text{(minimum)}, \qquad
  i(3{,}22) \approx 3{,}07\,\text{A} \ \text{(maximum)}.
\]
\end{exempel}
